\documentclass[a4paper,11pt]{article}

\usepackage{listings,xcolor}
 \lstset{language=Python, frame=tb}

\makeatletter
\let\@makecaption@ORI\@makecaption
\let\lst@MakeCaption@ORI\lst@MakeCaption
\makeatother

\usepackage[debug]{caption}

\makeatletter
\AtBeginDocument{%
\let\@makecaption@CAPTION\@makecaption
\let\lst@MakeCaption@CAPTION\lst@MakeCaption}
\newcommand\Original{\let\@makecaption\@makecaption@ORI\let\lst@MakeCaption\lst@MakeCaption@ORI}
\newcommand\Caption{\let\@makecaption\@makecaption@CAPTION\let\lst@MakeCaption\lst@MakeCaption@CAPTION}
\makeatother

\begin{document}

\section*{Regular caption}

\Original

\begin{lstlisting}[caption=Some caption,captionpos=t]
print()
\end{lstlisting}

\begin{lstlisting}[caption=Some caption,captionpos=b]
print()
\end{lstlisting}

\Caption\captionsetup{box=none}

\begin{lstlisting}[caption=Some caption,captionpos=t]
print()
\end{lstlisting}

\begin{lstlisting}[caption=Some caption,captionpos=b]
print()
\end{lstlisting}

\captionsetup{box=colorbox,boxcolor=orange!20}

\begin{lstlisting}[caption=Some caption,captionpos=t]
print()
\end{lstlisting}

\begin{lstlisting}[caption=Some caption,captionpos=b]
print()
\end{lstlisting}

\clearpage
\section*{float}

\Original

\begin{lstlisting}[float,caption=Eine Caption,frame=trbl]
Programmtext
\end{lstlisting}

\Caption\captionsetup{box=none}

\begin{lstlisting}[float,caption=Eine Caption,frame=trbl]
Programmtext
\end{lstlisting}

\captionsetup{box=colorbox,boxcolor=orange!20}

\begin{lstlisting}[float,caption=Eine Caption,frame=trbl]
Programmtext
\end{lstlisting}

\clearpage
\section*{linewidth=300pt}

\Original

\begin{lstlisting}[linewidth=300pt,caption=Some caption,captionpos=t]
print()
\end{lstlisting}

\begin{lstlisting}[linewidth=300pt,caption=Some caption,captionpos=b]
print()
\end{lstlisting}

\Caption\captionsetup{box=none}

\begin{lstlisting}[linewidth=300pt,caption=Some caption,captionpos=t]
print()
\end{lstlisting}

\begin{lstlisting}[linewidth=300pt,caption=Some caption,captionpos=b]
print()
\end{lstlisting}

\captionsetup{box=colorbox,boxcolor=orange!20}

\begin{lstlisting}[linewidth=300pt,caption=Some caption,captionpos=t]
print()
\end{lstlisting}

\begin{lstlisting}[linewidth=300pt,caption=Some caption,captionpos=b]
print()
\end{lstlisting}

\clearpage
\section*{xleftmargin=30pt, xrightmargin=10pt}

\Original

\begin{lstlisting}[xleftmargin=30pt,xrightmargin=10pt,caption=Some caption,captionpos=t]
print()
\end{lstlisting}

\begin{lstlisting}[xleftmargin=30pt,xrightmargin=10pt,caption=Some caption,captionpos=b]
print()
\end{lstlisting}

\Caption\captionsetup{box=none}

\begin{lstlisting}[xleftmargin=30pt,xrightmargin=10pt,caption=Some caption,captionpos=t]
print()
\end{lstlisting}

\begin{lstlisting}[xleftmargin=30pt,xrightmargin=10pt,caption=Some caption,captionpos=b]
print()
\end{lstlisting}

\captionsetup{box=colorbox,boxcolor=orange!20}

\begin{lstlisting}[xleftmargin=30pt,xrightmargin=10pt,caption=Some caption,captionpos=t]
print()
\end{lstlisting}

\begin{lstlisting}[xleftmargin=30pt,xrightmargin=10pt,caption=Some caption,captionpos=b]
print()
\end{lstlisting}

\clearpage
\section*{Enumeration, resetmargins=false}

\Original

\begin{enumerate}
\item Here is the first listing
\begin{lstlisting}[xleftmargin=30pt,resetmargins=false,caption=Some caption,captionpos=t]
print()
\end{lstlisting}
\begin{lstlisting}[xleftmargin=30pt,resetmargins=false,caption=Some caption,captionpos=b]
print()
\end{lstlisting}

\item And now a second one without \texttt{xleftmargin}
\begin{lstlisting}[xleftmargin=0pt,resetmargins=false,caption=Some caption,captionpos=t]
print()
\end{lstlisting}
\begin{lstlisting}[xleftmargin=0pt,resetmargins=false,caption=Some caption,captionpos=b]
print()
\end{lstlisting}
\end{enumerate}

\Caption\captionsetup{box=none}

\begin{enumerate}
\item Here is the first listing
\begin{lstlisting}[xleftmargin=30pt,resetmargins=false,caption=Some caption,captionpos=t]
print()
\end{lstlisting}
\begin{lstlisting}[xleftmargin=30pt,resetmargins=false,caption=Some caption,captionpos=b]
print()
\end{lstlisting}

\item And now a second one without \texttt{xleftmargin}
\begin{lstlisting}[xleftmargin=0pt,resetmargins=false,caption=Some caption,captionpos=t]
print()
\end{lstlisting}
\begin{lstlisting}[xleftmargin=0pt,resetmargins=false,caption=Some caption,captionpos=b]
print()
\end{lstlisting}
\end{enumerate}

\captionsetup{box=colorbox,boxcolor=orange!20}

\begin{enumerate}
\item Here is the first listing
\begin{lstlisting}[xleftmargin=30pt,resetmargins=false,caption=Some caption,captionpos=t]
print()
\end{lstlisting}
\begin{lstlisting}[xleftmargin=30pt,resetmargins=false,caption=Some caption,captionpos=b]
print()
\end{lstlisting}

\item And now a second one without \texttt{xleftmargin}
\begin{lstlisting}[xleftmargin=0pt,resetmargins=false,caption=Some caption,captionpos=t]
print()
\end{lstlisting}
\begin{lstlisting}[xleftmargin=0pt,resetmargins=false,caption=Some caption,captionpos=b]
print()
\end{lstlisting}
\end{enumerate}

\clearpage
\section*{Enumeration, resetmargins=true}

\Original

\begin{enumerate}
\item Here is the first listing
\begin{lstlisting}[xleftmargin=30pt,resetmargins=true,caption=Some caption,captionpos=t]
print()
\end{lstlisting}
\begin{lstlisting}[xleftmargin=30pt,resetmargins=true,caption=Some caption,captionpos=b]
print()
\end{lstlisting}

\item And now a second one without \texttt{xleftmargin}
\begin{lstlisting}[xleftmargin=0pt,resetmargins=true,caption=Some caption,captionpos=t]
print()
\end{lstlisting}
\begin{lstlisting}[xleftmargin=0pt,resetmargins=true,caption=Some caption,captionpos=b]
print()
\end{lstlisting}
\end{enumerate}

\Caption\captionsetup{box=none}

\begin{enumerate}
\item Here is the first listing
\begin{lstlisting}[xleftmargin=30pt,resetmargins=true,caption=Some caption,captionpos=t]
print()
\end{lstlisting}
\begin{lstlisting}[xleftmargin=30pt,resetmargins=true,caption=Some caption,captionpos=b]
print()
\end{lstlisting}

\item And now a second one without \texttt{xleftmargin}
\begin{lstlisting}[xleftmargin=0pt,resetmargins=true,caption=Some caption,captionpos=t]
print()
\end{lstlisting}
\begin{lstlisting}[xleftmargin=0pt,resetmargins=true,caption=Some caption,captionpos=b]
print()
\end{lstlisting}
\end{enumerate}

\captionsetup{box=colorbox,boxcolor=orange!20}

\begin{enumerate}
\item Here is the first listing
\begin{lstlisting}[xleftmargin=30pt,resetmargins=true,caption=Some caption,captionpos=t]
print()
\end{lstlisting}
\begin{lstlisting}[xleftmargin=30pt,resetmargins=true,caption=Some caption,captionpos=b]
print()
\end{lstlisting}

\item And now a second one without \texttt{xleftmargin}
\begin{lstlisting}[xleftmargin=0pt,resetmargins=true,caption=Some caption,captionpos=t]
print()
\end{lstlisting}
\begin{lstlisting}[xleftmargin=0pt,resetmargins=true,caption=Some caption,captionpos=b]
print()
\end{lstlisting}
\end{enumerate}

\clearpage
\section*{Description}

\lstset{%
        language=[Sharp]C,    % Sprache des Quellcodes ist Java
        numbers=left,
        stepnumber=1,            % Jede Zeile nummerieren.
        numbersep=5pt,          % 5pt Abstand zum Quellcode
        numberstyle=\tiny,      % Zeichengrösse 'tiny' für die Nummern.
        breaklines=true,        % Zeilen umbrechen wenn notwendig.
        breakautoindent=true,    % Nach dem Zeilenumbruch Zeile einrücken.
        numberblanklines=false,
        postbreak=\space,        % Bei Leerzeichen umbrechen.
        tabsize=2,              % Tabulatorgrösse 2        
        basicstyle=\ttfamily\scriptsize, % Nichtproportionale Schrift, klein für den Quellcode
        showspaces=false,        % Leerzeichen nicht anzeigen.
        showstringspaces=false,  % Leerzeichen auch in Strings ('') nicht anzeigen.
        extendedchars=true,      % Alle Zeichen vom Latin1 Zeichensatz anzeigen.
        framexleftmargin=5mm,
        frame=shadowbox,
        rulesepcolor=\color{red}
        }

\Original

\begin{description}
        \item[Test] Test text hier test text hiervest text hier test text hierest text hier test text hierest text hier test text hierest text hier test text
        
\begin{lstlisting}[caption=Titel des Codes, captionpos=t]
private void hideAllPlyBodies()
{
    foreach (HybridBody body in plyBodies)
    {
        hideElement(body, partDocument2);
    }

    foreach (HybridBody body in plyBodiesTriangularResults)
    {
        hideElement(body, partDocument2);
    }
}
\end{lstlisting}
\begin{lstlisting}[caption=Titel des Codes, captionpos=b]
private void hideAllPlyBodies()
{
    foreach (HybridBody body in plyBodies)
    {
        hideElement(body, partDocument2);
    }

    foreach (HybridBody body in plyBodiesTriangularResults)
    {
        hideElement(body, partDocument2);
    }
}
\end{lstlisting}

das ist der Text danach

\end{description}

\Caption\captionsetup{box=none}

\begin{description}
        \item[Test] Test text hier test text hiervest text hier test text hierest text hier test text hierest text hier test text hierest text hier test text
        
\begin{lstlisting}[caption=Titel des Codes, captionpos=t]
private void hideAllPlyBodies()
{
    foreach (HybridBody body in plyBodies)
    {
        hideElement(body, partDocument2);
    }

    foreach (HybridBody body in plyBodiesTriangularResults)
    {
        hideElement(body, partDocument2);
    }
}
\end{lstlisting}
\begin{lstlisting}[caption=Titel des Codes, captionpos=b]
private void hideAllPlyBodies()
{
    foreach (HybridBody body in plyBodies)
    {
        hideElement(body, partDocument2);
    }

    foreach (HybridBody body in plyBodiesTriangularResults)
    {
        hideElement(body, partDocument2);
    }
}
\end{lstlisting}

das ist der Text danach

\end{description}

\captionsetup{box=colorbox,boxcolor=orange!20}

\begin{description}
        \item[Test] Test text hier test text hiervest text hier test text hierest text hier test text hierest text hier test text hierest text hier test text
        
\begin{lstlisting}[caption=Titel des Codes, captionpos=t]
private void hideAllPlyBodies()
{
    foreach (HybridBody body in plyBodies)
    {
        hideElement(body, partDocument2);
    }

    foreach (HybridBody body in plyBodiesTriangularResults)
    {
        hideElement(body, partDocument2);
    }
}
\end{lstlisting}
\begin{lstlisting}[caption=Titel des Codes, captionpos=b]
private void hideAllPlyBodies()
{
    foreach (HybridBody body in plyBodies)
    {
        hideElement(body, partDocument2);
    }

    foreach (HybridBody body in plyBodiesTriangularResults)
    {
        hideElement(body, partDocument2);
    }
}
\end{lstlisting}

das ist der Text danach

\end{description}

\clearpage
\section*{Github issue \#1}

\lstset{%
        numbers=none,
        framexleftmargin=0pt}

\Original

\begin{lstlisting}[
	frame=shadowbox %helper, to see listing's dim
	,captionpos=b
	%,float %float is irrelevant to the issue
	%center the listing
		,xleftmargin=.20\textwidth
		,xrightmargin=.20\textwidth
%	,label=code
	,caption={This caption is not centered at all}
	,language=Delphi]
i := 1;
\end{lstlisting}

\Caption\captionsetup{box=none}

\begin{lstlisting}[
	frame=shadowbox %helper, to see listing's dim
	,captionpos=b
	%,float %float is irrelevant to the issue
	%center the listing
		,xleftmargin=.20\textwidth
		,xrightmargin=.20\textwidth
%	,label=code
	,caption={This caption is not centered at all}
	,language=Delphi]
i := 1;
\end{lstlisting}

\captionsetup{box=colorbox,boxcolor=orange!20}

\begin{lstlisting}[
	frame=shadowbox %helper, to see listing's dim
	,captionpos=b
	%,float %float is irrelevant to the issue
	%center the listing
		,xleftmargin=.20\textwidth
		,xrightmargin=.20\textwidth
%	,label=code
	,caption={This caption is not centered at all}
	,language=Delphi]
i := 1;
\end{lstlisting}

\clearpage
\section*{Gitlab issue \#26}

\lstset{frame=single}

\Original

\begin{lstlisting}[caption={Hello World!}]
int a = 1;
int b = 2;
\end{lstlisting}

\begin{lstlisting}[caption={Hello World!},captionpos=b]
int a = 1;
int b = 2;
\end{lstlisting}

\begin{description}
\item[A]
Lorem ipsum dolor sit amet, consectetur adipiscing elit, sed do eiusmod tempor incididunt ut labore et dolore magna aliqua.
\begin{lstlisting}[caption={Hello World!}]
int a = 1;
int b = 2;
\end{lstlisting}
Lorem ipsum dolor sit amet, consectetur adipiscing elit, sed do eiusmod tempor incididunt ut labore et dolore magna aliqua.
\begin{lstlisting}[caption={Hello World!},captionpos=b]
int a = 1;
int b = 2;
\end{lstlisting}
Lorem ipsum dolor sit amet, consectetur adipiscing elit, sed do eiusmod tempor incididunt ut labore et dolore magna aliqua.
\item[B]
Lorem ipsum dolor sit amet, consectetur adipiscing elit, sed do eiusmod tempor incididunt ut labore et dolore magna aliqua.
\end{description}

\clearpage
\Caption\captionsetup{box=none}

\begin{lstlisting}[caption={Hello World!}]
int a = 1;
int b = 2;
\end{lstlisting}

\begin{lstlisting}[caption={Hello World!},captionpos=b]
int a = 1;
int b = 2;
\end{lstlisting}

\begin{description}
\item[A]
Lorem ipsum dolor sit amet, consectetur adipiscing elit, sed do eiusmod tempor incididunt ut labore et dolore magna aliqua.
\begin{lstlisting}[caption={Hello World!}]
int a = 1;
int b = 2;
\end{lstlisting}
Lorem ipsum dolor sit amet, consectetur adipiscing elit, sed do eiusmod tempor incididunt ut labore et dolore magna aliqua.
\begin{lstlisting}[caption={Hello World!},captionpos=b]
int a = 1;
int b = 2;
\end{lstlisting}
Lorem ipsum dolor sit amet, consectetur adipiscing elit, sed do eiusmod tempor incididunt ut labore et dolore magna aliqua.
\item[B]
Lorem ipsum dolor sit amet, consectetur adipiscing elit, sed do eiusmod tempor incididunt ut labore et dolore magna aliqua.
\end{description}

\clearpage
\captionsetup{box=colorbox,boxcolor=orange!20}

\begin{lstlisting}[caption={Hello World!}]
int a = 1;
int b = 2;
\end{lstlisting}

\begin{lstlisting}[caption={Hello World!},captionpos=b]
int a = 1;
int b = 2;
\end{lstlisting}

\begin{description}
\item[A]
Lorem ipsum dolor sit amet, consectetur adipiscing elit, sed do eiusmod tempor incididunt ut labore et dolore magna aliqua.
\begin{lstlisting}[caption={Hello World!}]
int a = 1;
int b = 2;
\end{lstlisting}
Lorem ipsum dolor sit amet, consectetur adipiscing elit, sed do eiusmod tempor incididunt ut labore et dolore magna aliqua.
\begin{lstlisting}[caption={Hello World!},captionpos=b]
int a = 1;
int b = 2;
\end{lstlisting}
Lorem ipsum dolor sit amet, consectetur adipiscing elit, sed do eiusmod tempor incididunt ut labore et dolore magna aliqua.
\item[B]
Lorem ipsum dolor sit amet, consectetur adipiscing elit, sed do eiusmod tempor incididunt ut labore et dolore magna aliqua.
\end{description}

\end{document}

